\chapter{Open Problems}\label{ch:13}

In his 1900 address Hilbert remarked that a branch of science is alive as long as it offers an abundance of problems. Langlands' letter to Weil of 1967 and his 1970 paper \cite{langlands1970} supplied such an abundance, and more than half a century of work has settled only part of it. This chapter collects twelve problems on which the probabilistic point of view of this book has something definite to say. We do not pretend that probability will settle the central conjectures: the limits proposition (Proposition~\ref{prop:01limits}) shows that no functional of the empirical law of Frobenius traces can detect automorphy. What probability does supply is a language for quantitative refinements: rates, fluctuations, large deviations, extreme values, hitting times and sequential tests. For each problem we give the statement, its status with attributions, what probability and our computations contribute, and what would be needed. Where we do not know the right answer, the problem is posed as a question. Chapter~\ref{ch:14} collects computational versions of several of them. Chapter~\ref{ch:squarefree} poses two further problems, on the squarefree error under the Riemann hypothesis and on the signs of the P\'olya kernels of imaginary quadratic fields (Open Problems~\ref{op:sfquarter} and \ref{op:sfdensity}), Chapter~\ref{ch:beta} four on Dirichlet's $\beta$ function and the lines to the right of the centre (Open Problems~\ref{op:betacoupling} to \ref{op:betarecip}), and Chapter~\ref{ch:scholze} two on P\'olya kernels and signed Thorin measures (Open Problems~\ref{op:scholzekernel} and \ref{op:scholzelobe}). Chapter~\ref{ch:recip} poses eight problems on reciprocity (Open Problems~\ref{op:recipgln} to \ref{op:recipdrift}). Chapter~\ref{ch:hit} poses two on isogeny graphs and hitting times (Open Problems~\ref{op:hitsandpile} and \ref{op:hitgross}). Chapter~\ref{ch:flt} poses three, on the Hasse Weil bound, on patching and on Mazur's theorem (Open Problems~\ref{op:flthasse} to \ref{op:fltmazur}). Chapter~\ref{ch:lap} poses two, on the stable recovery of Frobenius traces from a window of the transported density and on penalised estimators of Thorin densities (Open Problems~\ref{op:lapwindow} and \ref{op:lapgg}). Chapter~\ref{ch:pd} poses three, on an upper bound for the rate of the Dirichlet process limit, on the signs of the Mertens deficits of fields with group $S_n$ and on an arithmetic counterpart of Bondesson's product theorem (Open Problems~\ref{op:pdrate} to \ref{op:pdproduct}). Chapter~\ref{ch:gl1} poses three, on the discrepancy of Hecke angles, on the Euler Kronecker constants of Hecke characters and on an effective Gamma shape test (Open Problems~\ref{op:gl1disc} to \ref{op:gl1cm}), and Chapter~\ref{ch:maass} four, on the Thorin representation at the centre, on a Chebotarev theorem for geodesics, on the even icosahedral representations and on the cutoff of Brownian motion on $X(N)$ (Open Problems~\ref{op:mscentre} to \ref{op:mscutoff}).

\section{Spectral problems}

\begin{problem}[Ramanujan Petersson for Maass forms]\label{op:rp}
Let $\phi$ be a Hecke Maass cusp form for $\SL_2(\Z)$ or for a congruence subgroup, with Satake parameters $\{\alpha_p,\alpha_p^{-1}\}$ at an unramified prime $p$, real Hecke eigenvalue $\lambda_\phi(p)=\alpha_p+\alpha_p^{-1}$ and Laplace eigenvalue $\lambda=\tfrac14+r^2$. Show that $|\alpha_p|=1$ for every such $p$, equivalently $|\lambda_\phi(p)|\le 2$, and that $r$ is real, which is Selberg's eigenvalue conjecture $\lambda\ge \tfrac14$.
\end{problem}

\paragraph{Status.} For holomorphic forms the analogue is Deligne's theorem \cite{deligne1974}; Maass forms have no motive and no geometric route is known.\index{Ramanujan Petersson conjecture} Selberg proved $\lambda\ge \tfrac{3}{16}$ and conjectured $\lambda\ge \tfrac14$ \cite{selberg1965}; for $\SL_2(\Z)$ itself the eigenvalue part is classical, since the cuspidal spectrum starts far above $\tfrac14$.\index{Selberg eigenvalue conjecture} The best general bounds, $|\alpha_p|\le p^{7/64}$ and $\lambda\ge \tfrac{975}{4096}$, are due to Kim and Sarnak \cite{kimsarnak2003}, from the automorphy of $\Sym^4$. Sarnak's density hypothesis \cite{sarnakxue1991}, proved in many cases, says that exceptional parameters are rare in a quantitative sense.

\paragraph{What probability and computation contribute.} By the trace formula the eigenvalues $\lambda_\phi(p)$, for $\phi$ on $\SL_2(\Z)$ ordered by Laplace eigenvalue, equidistribute with respect to the Plancherel measure $\mu_p$ \cite{sarnak1987}: the spectral measure at the root $o$ of simple random walk $S$ on the $(p+1)$ regular tree, rescaled to $[-2,2]$ (Chapter~\ref{ch:07}). Its moments $\int x^{2m}\,d\mu_p=(p+1)^{2m}p^{-m}\Prob(S_{2m}=o)$ grow like $4^m$ up to a polynomial factor. Hence if, in a family of $V$ forms, the averages of $\lambda_\phi(p)^{2m}$ match these moments up to $m\approx c\log V$, Markov's inequality bounds $\#\{\phi:|\lambda_\phi(p)|\ge2+\delta\}$ by $V(1+\delta/2)^{-2m}$ times a polynomial in $m$, that is, by $V^{1-2c\log(1+\delta/2)}$ up to such factors: the density hypothesis read as a large deviation bound for a random walk spectral measure. For the holomorphic analogue $\tau$, the primes below $10^6$ give $\max|\tau(p)/p^{11/2}|=1.99920$.

\paragraph{What would be needed.} Family averages cannot exclude a single exception. Individual bounds follow from automorphy of all $\Sym^n\phi$ by the tensor power trick \cite{langlands1970}; Newton and Thorne proved this automorphy for holomorphic newforms \cite{newtonthorne2021,newtonthorne2021b} through their Galois representations, which Maass forms in general lack.

\begin{problem}[Lang and Trotter]\label{op:lt}
Let $E/\Q$ be an elliptic curve of conductor $N_E$ without complex multiplication, let $t\in\Z$, and let $\pi_{E,t}(x)=\#\{p\le x:\ p\nmid N_E,\ a_p(E)=t\}$. Show that
\[
\pi_{E,t}(x)\ \sim\ C_{E,t}\,\pi_{1/2}(x),\qquad \pi_{1/2}(x)=\int_2^x\frac{du}{2\sqrt u\,\log u},
\]
whenever the constant $C_{E,t}$ predicted by Lang and Trotter \cite{langtrotter1976} is positive; when $C_{E,t}=0$ a congruence obstruction on the image of Galois makes the set finite.\index{Lang Trotter conjecture}
\end{problem}

\paragraph{Status.} Open for every pair $(E,t)$ with $C_{E,t}>0$. Elkies proved that every elliptic curve over $\Q$ has infinitely many supersingular primes \cite{elkies1987}, and, with Kaneko and M. R. Murty, that $\pi_{E,0}(x)=O(x^{3/4})$ \cite{elkies1991}; for every $t$, $\pi_{E,t}(x)=O(x(\log\log x)^2/(\log x)^2)$, as claimed by V. K. Murty \cite{murty1997} and proved by Thorner and Zaman \cite{thornerzaman2018}. Under the generalized Riemann hypothesis Zywina \cite{zywina2015}, improving Serre \cite{serre1981} and Murty, Murty and Saradha \cite{murtymurtysaradha1988}, proved $\pi_{E,0}(x)=O(x^{3/4}(\log x)^{-1/2})$ and $\pi_{E,t}(x)=O(x^{4/5}(\log x)^{-3/5})$. All are far above the predicted order $x^{1/2}/\log x$. On average over families the conjecture holds: Fouvry and Murty \cite{fouvrymurty1996} for $t=0$, David and Pappalardi \cite{davidpappalardi1999} for $t\ne0$.

\paragraph{What probability and computation contribute.} The conjecture is a local limit theorem with arithmetic corrections: the Sato Tate density at $0$ is $1/\pi$, the lattice spacing of $x_p$ is $1/\sqrt p$, and summing $1/(\pi\sqrt p)$ over primes gives $(2/\pi)\pi_{1/2}(x)$, which a Haar measure computation on the image of Galois in $\GL_2(\hat\Z)$ then corrects. For the curve $E_1\colon y^2=x^3+x+1$ of Chapter~\ref{ch:05} the counts of supersingular primes are $10$, $23$, $79$, $210$ at $x=10^4,\dots,10^7$, with ratios to $\pi_{1/2}(x)$ of $0.66$, $0.65$, $0.89$, $0.91$.

\paragraph{What would be needed.} Apart from the unconditional bound for $t=0$, which comes from singular moduli, the upper bounds above come from forms of the Chebotarev theorem in the division fields $\Q(E[\ell])$, whose error terms grow with the degree of the field even under the generalized Riemann hypothesis. An asymptotic requires control of a set of primes of density zero at the scale $x^{1/2}/\log x$, finer than any form of the Chebotarev theorem now available.

\begin{problem}[Optimal discrepancy and a functional limit theorem]\label{op:disc}
Let $E/\Q$ be without complex multiplication, let $p_1<p_2<\cdots$ be its primes of good reduction, $F_n$ the empirical distribution function of $x_{p_1},\dots,x_{p_n}$, $F_{\mathrm{ST}}$ that of $\ST$, and $D_n=\sup_x|F_n(x)-F_{\mathrm{ST}}(x)|$.
\begin{enumerate}[label=(\alph*)]
\item Show that $D_n=O(n^{-1/2+\delta})$ for every $\delta>0$.
\item (Question.) Let $E$ be chosen uniformly among the curves without complex multiplication of conductor at most $X$, and let $n=\lfloor X^c\rfloor$ with $c>0$ fixed. For which $c$ does $\sqrt n\,(F_n-F_{\mathrm{ST}})\circ F_{\mathrm{ST}}^{-1}$ converge in law to a Brownian bridge as $X\to\infty$, and $\sqrt n\,D_n$ to the Kolmogorov law?
\end{enumerate}
\end{problem}

\paragraph{Status.} Akiyama and Tanigawa \cite{akiyamatanigawa1999} conjectured (a) and showed that it implies the Riemann hypothesis for $L(s,E)$; Mazur \cite{mazur2008} stated the extension to every $L(s,\Sym^mE)$, under analytic hypotheses that Newton and Thorne now supply (Chapter~\ref{ch:05}). Unconditional effective forms of the Sato Tate theorem give only rates of logarithmic type \cite{thorner2021}, and the Riemann hypothesis for every $L(s,\Sym^mE)$ gives $D_n\ll n^{-1/4}$ up to logarithms (V. K. Murty \cite{murty1985}; Rouse and Thorner \cite{rousethorner2017} give an explicit form for squarefree level). For a single curve, with $n$ random on a logarithmic scale, the analogue of (b) has a negative answer under the Riemann hypothesis for $L(s,E)$, as explained below; the family version is, as far as we know, new.\index{discrepancy}\index{Brownian bridge}

\paragraph{What probability and computation contribute.} For independent draws (a) holds almost surely and (b) is Donsker's theorem. For a single curve, under the Riemann hypothesis for $L(s,E)$, the explicit formula argument of Section~\ref{sec:05disc} makes the first Chebyshev mode $n^{-1/2}\sum_{j\le n}x_{p_j}$ of order $(\log n)^{-1/2}$ in logarithmic mean square, whereas a Brownian bridge limit would make this continuous linear functional of the process standard normal; hence no randomisation of $n$ on a logarithmic scale, such as $n$ logarithmically uniform on $[N,N^2]$, produces a Brownian bridge for one curve; under the Riemann hypothesis for $L(s,\Sym^mE)$ every fixed mode $n^{-1/2}\sum_{j\le n}U_m(x_{p_j}/2)$ tends to $0$ in the same way, and Conjecture~\ref{conj:05bridge} predicts that $\sqrt n\,D_n$ itself tends to $0$ in logarithmic density. In a family one expects the primes up to $p_n$ to behave independently while $p_n$ is small compared with the conductors; a central limit theorem for the first mode is known for eigenforms of level one whose weight grows faster than every power of $n$ \cite{nagoshi2006}. The explicit formula takes over once $p_n$ exceeds the conductor of $\Sym^mE$, about $N_E^m$, and its diagonal terms suggest that the $m$th mode has variance about $\min(1,m/c)$: a transition at $c=1$, below which every mode has variance near $1$ and above which the first mode is too small for a bridge. The data of Section~\ref{sec:05disc} are consistent with so slow a decay: $\sqrt n\,D_n$ shows no drift for $E_1$, and across $200$ curves a Kolmogorov Smirnov test does not separate its law from the Kolmogorov law.

\paragraph{What would be needed.} For (a), square root cancellation in the mode sums $\sum_{j\le n}U_m(x_{p_j}/2)$ uniformly for $m$ up to about $n^{1/2}$; the Riemann hypothesis for each $L(s,\Sym^mE)$ bounds the $m$th mode sum only by $m\sqrt n$ up to logarithms, because the number of zeros of bounded height grows linearly in $m$, and this is why it yields only $n^{-1/4}$. For (b), in the regime $c<1$ a central limit theorem for the modes jointly as $m$ grows, with tightness; in the regime $c>1$ the variances of the low modes, which depend on the pair correlation of the low lying zeros of the symmetric power $L$ functions in the family.

\section{Inverse problems and symmetry}

\begin{problem}[An inverse problem for automorphy]\label{op:inverse}
(Question.) Is there a criterion, stated through quantitative properties of the twisted sums $\sum_{p\le x}x_p\chi(p)\log p$ at the square root scale (rates in the law of large numbers, mean squares on the logarithmic scale, large deviation rates), uniform over Dirichlet characters $\chi$ of growing conductor, that decides whether a sequence $(x_p)$ in $[-2,2]$ indexed by the primes is close to the normalised Hecke eigenvalues $x_p(f)=a_p(f)/p^{(k-1)/2}$ of a holomorphic newform $f$ of weight $k$ with trivial nebentypus, in the sense that $\sum_{p\le x}(x_p-x_p(f))\chi(p)\log p=o(\sqrt x)$ for every $\chi$?
\end{problem}

\paragraph{Status.} The converse theorems of Hecke \cite{hecke1936}, Weil \cite{weil1967}, Jacquet and Langlands \cite{jacquetlanglands1970} and Cogdell and Piatetski Shapiro \cite{cogdellps1994} are inverse theorems of this type, but analytic: they start from the analytic continuation and functional equations of twisted $L$ functions. Chapter~\ref{ch:recip} restates the hypotheses of Hecke's theorem as a symmetry of a P\'olya kernel built from the coefficients (Proposition~\ref{prop:recipsym}), and the holomorphy and the generalized Riemann hypothesis of a self dual Artin $L$ function as the complete monotonicity of a heat trace computed from Frobenius traces, the conductor and the gamma factor (Theorem~\ref{thm:recipcm}); these are reformulations, and Open Problems~\ref{op:recipgln} and \ref{op:reciplaw} pose related questions. Chapter~\ref{ch:lap} computes the P\'olya kernel of an Artin $L$ function by Bromwich inversion on the critical line, given that the $L$ function is entire; that the result is modular is again the content of Weil's converse theorem (Proposition~\ref{prop:lappolya}).\index{converse theorem}

\paragraph{What probability and computation contribute.} By Corollary~\ref{cor:01functional} and the remark after it, no statistic that is unchanged when the $x_p$ are altered on a set of primes of density zero, such as a limit of averages with weight $O(1/\pi(x))$ per prime, twisted or not, can decide automorphy. Rates at the square root scale cannot decide it either, which is why the question asks only for closeness: if $S$ is an infinite set of primes with $\theta_S(x)=\sum_{p\in S,\,p\le x}\log p=o(\sqrt x)$, such as the primes $p_{2^j}$, altering the $x_p(f)$ within $[-2,2]$ at every prime of $S$ changes each twisted sum by at most $4\theta_S(x)=o(\sqrt x)$ and leaves every such rate unchanged, while by Proposition~\ref{prop:01limits} the altered sequence agrees at all but finitely many primes with no such newform. Pointwise bounds do not suffice: independent variables with the Sato Tate law satisfy, almost surely, $\sum_{p\le x}x_p\chi(p)\log p\ll\sqrt x\log(qx)$ for all characters $\chi$ of conductor $q\le x$, and are almost surely close to no newform, since the newforms are countable and for each of them the untwisted sum of the differences fluctuates on the scale $\sqrt{x\log x}$. Mean squares on the logarithmic scale separate them: under the generalized Riemann hypothesis, for a newform $f$ the function $e^{-y/2}\sum_{p\le e^y}x_p(f)\chi(p)\log p$ is almost periodic in mean square, with frequencies the ordinates of the zeros of $L(s,f\otimes\chi)$, and so bounded in mean square, while for independent variables its mean square grows like $y$. Square root cancellation for every twist, even with a logarithmic loss, makes the twisted $L$ functions holomorphic and nonvanishing in $\operatorname{Re}s>\tfrac12$, but it does not give their functional equations.

\paragraph{What would be needed.} A theorem converting such rates into the hypotheses of Weil's converse theorem, not for $(x_p)$ itself, which may carry an alteration that the rates cannot see, but for a sequence reconstructed from the rates, for instance from the frequencies above; we do not know whether such a statement can be true without the root numbers.

\begin{problem}[General Sato Tate groups and symmetry types]\label{op:stgroups}
(a) For abelian varieties and motives over number fields, determine the Sato Tate group and prove equidistribution with respect to its Haar measure. (b) For natural families of $L$ functions, prove the Katz and Sarnak predictions for low lying zeros for test functions of unrestricted Fourier support.
\end{problem}

\paragraph{Status.} For elliptic curves without complex multiplication over totally real fields, (a) is the theorem of Clozel, Harris, Shepherd Barron and Taylor, completed by Barnet Lamb, Geraghty, Harris and Taylor \cite{cht2008,hsbt2010,taylor2008,blght2011}; over CM fields it was proved by Allen, Calegari, Caraiani and seven coauthors \cite{allen2023}. For abelian surfaces there are exactly $52$ groups \cite{fkrs2012}; equidistribution is known for many with extra endomorphisms \cite{johansson2017} but not for $\USp(4)$. For abelian threefolds Fit\'e, Kedlaya and Sutherland announced a classification with $410$ groups \cite{fks2021}. For (b), Katz and Sarnak \cite{katzsarnak1999,katzsarnak1999b} prove the function field analogue; over number fields only restricted support is known \cite{ils2000}, as for the $n$ level correlations of a single $L$ function \cite{rudnicksarnak1996}.\index{Sato Tate group}\index{symmetry type}

\paragraph{What probability and computation contribute.} Both parts concern pushforwards of Haar measure: in (a) to the conjugacy classes of the group, in (b) to the eigenvalue statistics of the symmetry group. In (b), restricted support means that the $n$ level densities are known only against test functions whose Fourier transforms have bounded support, which falls short of determining them, much as finitely many moments fall short of determining a law. The Weyl averages of $E_1$ and of $E_2\colon y^2=x^3-x$ in Figure~\ref{fig:04weyl} of Chapter~\ref{ch:04} match the generic and the CM law; how fast such agreement identifies the group is Open Problem~\ref{op:bayes}.

\paragraph{What would be needed.} For (a), already for a generic abelian surface, potential automorphy of the Galois representation composed with every irreducible representation of $\USp(4)$, not only the symmetric powers; for (b), a number field substitute for monodromy.

\begin{problem}[Functoriality beyond known cases]\label{op:funct}
Prove functoriality for $\GL_m\times\GL_n\to\GL_{mn}$ in general, and for $\Sym^n:\GL_2\to\GL_{n+1}$ for Maass forms and all $n$.
\end{problem}

\paragraph{Status.} Known cases include $\Sym^2$ \cite{gelbartjacquet1978}, $\GL_2\times\GL_2$ \cite{ramakrishnan2000}, $\GL_2\times\GL_3$ and $\Sym^3$ \cite{kimshahidi2002}, $\Sym^4$ (Kim \cite{kimsarnak2003}), all symmetric powers for holomorphic newforms \cite{newtonthorne2021,newtonthorne2021b}, and the transfers from quasi split classical groups to $\GL_N$ \cite{arthur2013}; see \cite{arthur2003,gelbart1984}.\index{functoriality}

\paragraph{What probability and computation contribute.} Functoriality has consequences that Frobenius data can test. For an elliptic curve without complex multiplication, automorphy and cuspidality of every symmetric power give the Sato Tate law, and the even moments of $x_p$ for $E_1$ over the good primes below $10^7$ agree with the Catalan numbers $1,2,5,14,42,132$ to a relative error below $6\times10^{-4}$ (Chapter~\ref{ch:05}). For cuspidal $\pi_1,\pi_2$ on $\GL_2$, cuspidality of $\pi_1\boxtimes\pi_2$ on $\GL_4$ gives
\[
\sum_{p\le x}\frac{|\tr A_p(\pi_1)|^2\,|\tr A_p(\pi_2)|^2}{p}\sim\log\log x,
\]
which for two elliptic curves without complex multiplication that are not twists of each other is the moment identity $\int\!\!\int x^2y^2\,d\ST(x)\,d\ST(y)=1$.

\paragraph{What would be needed.} The known cases came from the Langlands Shahidi method and converse theorems (Gelbart and Jacquet, Ramakrishnan, Kim and Shahidi, Kim), from the trace formula (Arthur) and from automorphy lifting (Newton and Thorne). By the limits proposition the checks above are functionals of empirical laws and cannot replace such methods.

\section{Graphs}

\begin{problem}[Ramanujan graphs of all degrees]\label{op:graphs}
Give a deterministic algorithm that, for each $d\ge 3$ and all large $n$ with $nd$ even, constructs in time polynomial in $n$ a nonbipartite $d$ regular graph on $n$ vertices whose nontrivial eigenvalues satisfy $|\lambda|\le 2\sqrt{d-1}$.
\end{problem}

\paragraph{Status.} The arithmetic constructions of Chapter~\ref{ch:08} (Lubotzky, Phillips and Sarnak; Margulis; Morgenstern; Pizer) give special sequences of $n$ only; the probabilistic results of that chapter give nearly Ramanujan graphs (Friedman, Bordenave), bipartite Ramanujan multigraphs of every degree and every even number of vertices (Marcus, Spielman and Srivastava) and, since about $69$ percent of random $d$ regular graphs are Ramanujan (Huang, McKenzie and Yau~\cite{hmy2024}, a 2024 preprint), existence for all large $n$ and a randomised polynomial time construction (sample a graph and test its spectrum). Deterministic polynomial time constructions reach $|\lambda|\le2\sqrt{d-1}+\varepsilon$ for every $d\ge3$ and $\varepsilon>0$ (Mohanty, O'Donnell and Paredes \cite{mop2022}), on exactly $n$ vertices for all large $n$ with $nd$ even (Alon \cite{alon2021}); the problem asks for $\varepsilon=0$, which, as far as we know, no deterministic construction achieves for nonbipartite graphs of all large sizes.\index{Ramanujan graph}

\paragraph{What probability and computation contribute.} The existence results beyond the arithmetic constructions are probabilistic: the trace method of Friedman and Bordenave, the interlacing polynomials of Marcus, Spielman and Srivastava, and the edge universality around the Kesten McKay law of Huang, McKenzie and Yau.

\paragraph{What would be needed.} Derandomisation of the random constructions, which Cohen \cite{cohen2016} achieved for the bipartite construction of Marcus, Spielman and Srivastava, or an arithmetic source of graphs of every degree: the arithmetic constructions above come from groups acting on the tree of $\PGL_2$ over a local field, whose degree is always a prime power plus one.

\section{Thorin measures and extreme values}

\begin{problem}[Hilbert's eighth problem for $\GL_2$ in Thorin form]\label{op:thorin}
Let $\Lambda(s,\Delta)=(2\pi)^{-s}\Gamma(s)L(s,\Delta)$, write its zeros as $6\pm i\gamma_k$ with $\operatorname{Re}\gamma_k>0$, and let $k_\Delta(x)=\sum_k e^{-\gamma_k^2x}$ for $x>0$. Prove by a probabilistic argument that $k_\Delta$ is completely monotone on $(0,\infty)$.
\end{problem}

\paragraph{Status.} Integrating termwise, $\sum_k(\gamma_k^2+s)^{-1}=\int_0^\infty e^{-sx}k_\Delta(x)\,dx$ for $s>0$, and by the Hadamard product the left side is $(d/ds)\log\Lambda(6+\sqrt s,\Delta)$. By Bernstein's theorem and the equivalence with generalized gamma convolutions in Chapter~\ref{ch:10} (Theorem~\ref{thm:ch10ggc}) the problem is equivalent to the Riemann hypothesis for $\Delta$, hence open. The parallel statement for $\xi$ is Polson's reformulation of Hilbert's eighth problem \cite{polson2017}, building on Biane, Pitman and Yor \cite{bpy2001}; the theory of generalized gamma convolutions is due to Thorin \cite{thorin1977} and Bondesson \cite{bondesson1992}.\index{Thorin measure}\index{complete monotonicity}

\paragraph{What probability and computation contribute.} A probabilistic proof would exhibit the law with Laplace transform $\Lambda(6)/\Lambda(6+\sqrt s)$, whose Lévy density is $x^{-1}k_\Delta(x)$, as a generalized gamma convolution built without the zeros, for instance from hitting or occupation times of a diffusion, as Bessel processes produce the laws of Biane, Pitman and Yor. Equivalently, the characteristic function of the Pólya variable $W$ must have only real zeros, which is Pólya's question \cite{polya1926}, with the de Bruijn Newman heat flow \cite{debruijn1950,newman1976,rodgerstao2020} as a natural deformation. The product over the $404$ zeros up to height $400$, with a tail correction, reproduces $\Lambda(6)/\Lambda(6+\sqrt s)$ for $s\le400$ to about four significant digits (relative error at most $1.9\times10^{-4}$), and $\operatorname{Var}W=0.112352$ against $0.112348$ from the zeros with the tail correction (Chapter~\ref{ch:10}). Chapter~\ref{ch:scholze} computes a smoothed signed version of this Thorin measure by Stieltjes inversion to height $30$.

\paragraph{What would be needed.} A positive measure $\tau$ on $(0,\infty)$ with $k_\Delta(x)=\int e^{-tx}\,\tau(dt)$, the Thorin measure of the problem, constructed without reference to the zeros, for example from an occupation or hitting time law; such a construction would prove the Riemann hypothesis for $\Delta$. Positivity of $k_\Delta$ alone, which is infinite divisibility, would not suffice. Open Problem~\ref{op:recipfejer} asks whether a positive summation of the prime side of the explicit formula can play the part that Fej\'er's kernel plays for the circle (Proposition~\ref{prop:recipfejer}).

\begin{problem}[The extreme value representation]\label{op:ev}
Let $M^*$ be the maximum of $24$ independent copies of $M=\max_{n\ge1}E_n/(2\pi n)$, with $E_n$ independent standard exponential variables, so that $L(s,\Delta)=\E\,Q(s,2\pi M^*)$ with $Q(s,c)=\Gamma(s,c)/\Gamma(s)$ (Chapter~\ref{ch:11}).
\begin{enumerate}[label=(\alph*)]
\item (Question.) Can properties of the law of $M^*$, such as the self similarity $F(y)=y^{-12}e^{2\pi(y-1/y)}F(1/y)$ of its distribution function given by modularity, be used to locate or count the zeros of $L(s,\Delta)$?
\item (Question.) For which newforms $f$ is there a random variable $M_f$, the maximum of countably many independent variables with distribution functions of the form $(1-e^{-cy})^{\kappa}$ with $c,\kappa>0$, such that $f(iy)=e^{-2\pi y}\Prob(M_f\le y)$ for all $y>0$?
\end{enumerate}
\end{problem}

\paragraph{Status.} The representation is proved in Chapter~\ref{ch:11} for $\Delta$ and for every eta product $\prod_m\eta(m\tau)^{r_m}$ with $r_m\ge0$ and $\sum_m m\,r_m=24$, among them the multiplicative eta products of Dummit, Kisilevsky and McKay \cite{dkm1985} (Theorems~\ref{thm:ch11main} and \ref{thm:ch11general}), so (b) has a positive answer for those of them that are newforms. Part (a) asks for a route to Conjecture~\ref{conj:ch11rh}. As far as we know, both questions are new.\index{extreme value representation}

\paragraph{What probability and computation contribute.} By Corollaries~\ref{cor:ch11bounds} and~\ref{cor:ch11monotone}, $L(s,\Delta)$ lies in $(0,1)$ and increases on $s>0$. On the critical line the functional equation gives
\[
\Lambda(6+it,\Delta)=2\,\E\int_{\max(1,M^*)}^\infty e^{-2\pi y}y^5\cos(t\log y)\,dy,
\]
so (a) asks for the sign changes in $t$ of an explicit expectation, the analogue of Hardy's function. By uniqueness of Laplace transforms, (b) holds exactly when the exponents $\kappa_m$ in the formal product $q^{-1}f=\prod_{m\ge1}(1-q^m)^{\kappa_m}$ are all nonnegative and the product converges for $0<q<1$; for nonnegative $\kappa_m$ this convergence is equivalent to $f$ having no zeros in the upper half plane. Chapter~\ref{ch:11} checks the representation to $16$ digits and by Monte Carlo.

\paragraph{What would be needed.} For (a), a property of the specific law of $M^*$ that fails for nearby laws, since expectations built from nearby laws need not have their zeros on a line. For (b), a description of the newforms with no zeros in the upper half plane, whose divisors are supported on the cusps, together with the signs of their exponents $\kappa_m$; eta quotients are the basic examples, Borcherds products are a natural framework for the exponents, and Martin's classification of the multiplicative eta quotients \cite{martin1996} is the place to start.

\section{Hitting times, inference and bias}

\begin{problem}[Kemeny constants of isogeny graphs at finite level]\label{op:kemeny}
For primes $N\equiv1\bmod 12$ let $K_N$ be the Kemeny constant of simple random walk on the supersingular $2$ isogeny graph in characteristic $N$, which is $3$ regular on $h=(N-1)/12$ vertices, so that $K_N/h\to2$ (Theorem~\ref{thm:08kemeny}). (Question.) Does the empirical distribution of $K_N-2h$ over the primes $N\le Y$ in this class converge as $Y\to\infty$, and is $K_N-2h$ bounded?
\end{problem}

\paragraph{Status.} The first order law is Theorem~\ref{thm:08kemeny} of Chapter~\ref{ch:08}, proved from Serre's vertical equidistribution \cite{serre1997} and the Brandt matrix identification \cite{pizer1990}. The second order term has not been studied as far as we know; hitting times matter for isogeny based hash functions \cite{cgl2009}.\index{Kemeny constant}

\paragraph{What probability and computation contribute.} By Theorem~\ref{thm:A:kemeny} and Remark~\ref{rem:08exact},
\[
K_N=\sum_{j\ge2}\frac{1}{1-\lambda_j/3}=\frac32\sum_{m\ge0}2^{-m}\tr\bigl(T_{2^m}\mid S_2(\Gamma_0(N))\bigr),
\]
where $\lambda_2,\dots,\lambda_h$ are the eigenvalues of $T_2$ on $S_2(\Gamma_0(N))$. The Eichler Selberg trace formula writes $\tr T_{2^m}$ as $(N+1)/12$ for even $m$, plus terms independent of $N$, minus half the sum of the Hurwitz class numbers $H(2^{m+2}-t^2)$ weighted by $1+\bigl(\tfrac DN\bigr)$, $D=t^2-2^{m+2}$. The main terms add up to $(N+1)/6=2h+\tfrac13$, so that, formally, $K_N-2h=c_0+\sum_Dw_D\bigl(\tfrac DN\bigr)$ with explicit constants $c_0$ and $w_D\le0$, the sum running over negative fundamental discriminants $D\ne-3,-4$. On this progression $\bigl(\tfrac DN\bigr)$ is the product of the signs $\varepsilon_q=\bigl(\tfrac qN\bigr)$ over the primes $q\ne3$ that divide $D$ to an odd power, and by Dirichlet's theorem, for each finite set of primes $q$, the signs $\varepsilon_q$ are asymptotically independent and fair as $N$ varies. This suggests that $K_N-2h$ converges in law to the random multiplicative series obtained by taking the $\varepsilon_q$ independent. Size estimates suggest that the weights are square summable, so that this series converges in mean square, but not absolutely summable, so that $K_N-2h$ may be unbounded, with large negative values at levels where $\varepsilon_q=1$ for many small primes $q$. Over the $207$ levels $1000<N<8000$ computed for this book, $K_N-2h$ lies between $-17.7$ and $6.9$ with mean $-7.0$; its mean is $-11.0$ when $\bigl(\tfrac{-7}{N}\bigr)=1$ and $-3.3$ when $\bigl(\tfrac{-7}{N}\bigr)=-1$ (Chapter~\ref{ch:08}).

\paragraph{What would be needed.} Control of the trace formula with weights $2^{-m}$ up to $m$ of order $2\log_2N$, beyond which the Ramanujan bound controls the tail; cancellation in the class number sums for discriminants larger than $N$, where the symbols $\bigl(\tfrac DN\bigr)$ are periodic in $D$ rather than independent; and joint equidistribution of the signs $\varepsilon_q$ for $q$ up to a power of $N$, a large sieve statement.

\begin{problem}[Bayesian identification of Sato Tate groups]\label{op:bayes}
(a) Extend Theorem~\ref{thm:12consist} to the $52$ Sato Tate groups of abelian surfaces, whose laws of the normalised coefficients $(a_1,a_2)$ can have atoms and components on lower dimensional sets, and to motives of higher rank. (b) Make it effective: bound the number of primes after which the log Bayes factor favours the true group against every other candidate, polynomially in the logarithm of the conductor and without assuming the Riemann hypothesis.
\end{problem}

\paragraph{Status.} Theorem~\ref{thm:12consist} gives the double limit $\lim_{M\to\infty}\lim_{n\to\infty}n^{-1}\log B_n^M=\mathrm{KL}(\mu_0\|P_0)-\mathrm{KL}(\mu_0\|P_1)$ for the log Bayes factor $\log B_n^M$ of $P_1$ against $P_0$ clipped at level $M$, when $P_0$, $P_1$ and the limit law $\mu_0$ of the data have densities on a compact set with respect to Lebesgue measure plus counting measure on a fixed finite set $D$ and the data are equidistributed for $\mu_0$ relative to $D$, under the further conditions stated there; Chapter~\ref{ch:12} applies it to elliptic curves and to three candidates in genus $2$. The classification \cite{fkrs2012} matches curves to groups by moments, not by sequential tests; several of the $52$ groups share the law of $a_1$, so the pair $(a_1,a_2)$ is needed. For independent observations Wald's test \cite{wald1947} is optimal for simple hypotheses \cite{waldwolfowitz1948}, with error exponents given by the divergences. For elliptic curves, effective forms of the Sato Tate theorem \cite{rousethorner2017,thorner2021} yield bounds of the kind asked in (b), polynomial in $\log N_E$ under the Riemann hypothesis for the symmetric power $L$ functions and much weaker unconditionally.\index{sequential probability ratio test}\index{Kullback Leibler divergence}

\paragraph{What probability and computation contribute.} With $\mu_0=P_0=\mu_G$ for the true group $G$ and $P_1=\mu_{G'}$ for a candidate $G'$, the double limit is $-d_{GG'}$ with $d_{GG'}=\mathrm{KL}(\mu_G\|\mu_{G'})$, so a sequential test with threshold $\log(1/\beta)$ stops after about $\log(1/\beta)/d_{GG'}$ primes. When $\mu_G$ charges a set that $\mu_{G'}$ does not, the divergence is infinite, and the small atoms and clipping of Chapter~\ref{ch:12} make it large and finite; the near ties, which set the error exponents, come from groups that differ only on components of small mass. The $52\times52$ divergence matrix is part of Computational Problem 29.2.

\paragraph{What would be needed.} For (a), reference measures adapted to the lower dimensional components, so that the clipping argument still applies; equidistribution for the true group, which is open for $\USp(4)$; and, in higher rank, the classification of the groups. For (b), unconditional effective equidistribution with power rates, which would also bear on Open Problem~\ref{op:disc}.

\begin{problem}[The rank bias]\label{op:rankbias}
Let $E/\Q$ have conductor $N_E$ and analytic rank $r_E$, let $S_E(x)=\sum_{p\le x,\,p\nmid N_E}a_p\log p/p$, and assume the Riemann hypothesis for $L(s,E)$. Then $S_E(x)-(\tfrac12-r_E)\log x$, with $\log x$ uniform on $[0,Y]$ and $Y\to\infty$, converges in law to a distribution $\mu_E$ with mean $C_E$; let $Z_E$ have the law $\mu_E$ shifted to mean zero. (Question.) As $E$ runs through the curves of conductor at most $X$ and analytic rank $r\in\{0,1\}$, does the law of $Z_E/\log N_E$ converge as $X\to\infty$, and is the limit the one predicted by the low lying zeros of the corresponding orthogonal family?
\end{problem}

\paragraph{Status.} Under the Riemann hypothesis for $L(s,E)$ the explicit formula and the method of Rubinstein and Sarnak \cite{rubinsteinsarnak1994} give the existence of $\mu_E$, and identify it as the law of $C_E-\sum_{\gamma>0}2\sin(\theta_\gamma)/\gamma$, with independent uniform phases $\theta_\gamma$, when the positive ordinates $\gamma$ of the noncentral zeros are linearly independent over $\Q$; Sarnak's letter \cite{sarnak2007} explains such biases through the explicit formula. The link between the rank and partial Euler products goes back to Goldfeld \cite{goldfeld1982} and Kuo and Murty \cite{kuomurty2005}; see Mazur \cite{mazur2008}. The family question is, as far as we know, open.\index{rank bias}\index{Chebyshev bias}

\paragraph{What probability and computation contribute.} The explicit formula gives $S_E(x)-(\tfrac12-r_E)\log x=C_E-\sum_{\gamma>0}2\sin(\gamma\log x)/\gamma+o(1)$ in logarithmic mean square: the prime squares produce the $\tfrac12$, the central zero of order $r_E$ the $-r_E$, and the other zeros an almost periodic oscillation (Chapter~\ref{ch:05}). Under linear independence $Z_E$ has variance $V_E=2\sum_{\gamma>0}\gamma^{-2}$, the same form as $\operatorname{Var}W$ in Open Problem~\ref{op:thorin}. For the four curves of Chapter~\ref{ch:05}, the central Taylor coefficients computed in Chapter~\ref{ch:10b} give $V_E=0.255$, $0.379$, $0.814$ and $1.452$ for 11a1, 37a1, 389a1 and 5077a1, increasing as the lowest noncentral zero descends from $6.36$ to $2.05$. The fitted slopes $0.46$, $-0.59$, $-1.43$, $-2.51$ against $0.5$, $-0.5$, $-1.5$, $-2.5$ (Chapter~\ref{ch:05}) are consistent with oscillations of standard deviation $\sqrt{V_E}$, between $0.5$ and $1.2$, over the range $6.9\le\log x\le16.1$. Since the lowest zeros of a curve of large conductor lie at heights of order $1/\log N_E$, $V_E$ is typically of order $(\log N_E)^2$, which explains the normalisation. If the normalised ordinates $t_j=\gamma_j\log N_E/(2\pi)$ follow the Katz Sarnak limit, the even orthogonal one for $r=0$ and the odd one with its central zero removed for $r=1$, then $Z_E/\log N_E$ should converge in law to $-\pi^{-1}\sum_j\sin(\theta_j)/t_j$, with $(t_j)$ the points of that limit and $\theta_j$ independent and uniform. For $r=0$ this limit has infinite variance, because the even orthogonal ensemble does not repel the lowest zero from the central point; for $r=1$ the repulsion makes the variance finite. The slope test at high rank is Computational Problem 5.3.

\paragraph{What would be needed.} Each $\mu_E$ needs the Riemann hypothesis for $L(s,E)$. For the family law, $V_E$ is dominated by the lowest zeros, so one needs the distribution of the lowest zero itself, which the restricted support results for the one level density \cite{ils2000} do not determine (Open Problem~\ref{op:stgroups}).

\section{Notes}

The conjectures behind Open Problems~\ref{op:inverse} and \ref{op:funct} go back to Langlands' paper \cite{langlands1970}; Arthur \cite{arthur2003} and Gelbart \cite{gelbart1984} are standard overviews. Sarnak's density hypothesis and vertical Plancherel law \cite{sarnak1987,sarnakxue1991} are the model for the probabilistic reading of Open Problem~\ref{op:rp}. Lang and Trotter \cite{langtrotter1976} derived their constant from a probabilistic model of Frobenius in the image of Galois, the model behind Open Problem~\ref{op:lt}. Mazur's essay \cite{mazur2008} is the source of the error term perspective behind Open Problems~\ref{op:disc} and \ref{op:rankbias}; in both, the explicit formula makes the arithmetic sums smaller and more structured than independent ones, the theme of the analysis of Chebyshev's bias by Rubinstein and Sarnak \cite{rubinsteinsarnak1994} and of Sarnak's letter \cite{sarnak2007}. The Thorin and extreme value problems are specific to this book.

\section{Open computational problems}

The full programme is in Chapter~\ref{ch:14}; the two problems below belong to this chapter.

\begin{compproblem}[Discrepancy at scale]\label{cp:13disc}
Objective: test the transition at $c=1$ suggested in Open Problem~\ref{op:disc}(b), and the slow decay predicted for single curves. It extends Computational Problem~\ref{cp:05bridge}, which runs $E_1$ and a family to a fixed prime bound, by tying the number of primes to the conductor, $n=\lfloor X^c\rfloor$. Scale: for $E_1$ and $10$ further curves without complex multiplication, all primes below $10^{10}$; for the family, $10^3$ curves with conductor in $[X,2X]$ from the LMFDB \cite{lmfdb} for $X=10^3$ and $10^4$, with $n=\lfloor X^c\rfloor$ for $c=\tfrac12,1,\tfrac32,2$, so at most $10^8$ primes per curve. Route: average polynomial time point counting \cite{harvey2014,harveysutherland2014}, or PARI/GP \texttt{ellap} in parallel. Deliverables: for each $X$ and $c$, the variances across the family of the modes $n^{-1/2}\sum_{j\le n}U_m(x_{p_j}/2)$ for $m\le20$, compared with $\min(1,m/c)$, and the law of $\sqrt n\,D_n$ against the Kolmogorov law with a Kolmogorov Smirnov statistic; for the single curves, the paths of $\sqrt n\,D_n$ and of the modes for $m\le200$ in $\log n$. Why it matters: the $200$ curves of Chapter~\ref{ch:05} have discriminants below $10^7$ in absolute value and use the primes below $10^6$, so $c$ is of order $1$ for them, and their standard deviation $0.219$ against $0.2603$ is consistent with a mild deficit there; varying $c$ separates the two regimes, and a deficit for $c>1$ would indicate that the explicit formula, not independence, governs the fluctuations.
\end{compproblem}

\begin{compproblem}[Second order Kemeny constants]\label{cp:13kemeny}
Objective: test the random series model of Open Problem~\ref{op:kemeny}, which Computational Problem~\ref{cp:08second} approaches through a regression on the symbols for small discriminants, whereas here the full predicted series is tested. Data: the values of $K_N-2h$ for the $19564$ primes $N\equiv1\bmod12$ below $10^6$ ($h$ up to $83330$) produced in Computational Problem~\ref{cp:08second}. Route: compute the weights $w_D$ for $|D|\le10^6$ from the trace formula, as $-\tfrac34$ times the sum of $2^{-m}H(2^{m+2}-t^2)$ over the pairs $(m,t)$ for which $t^2-2^{m+2}$ is $D$ times a square, with PARI/GP \texttt{qfbhclassno}; compare $K_N-2h$ with $c_0+\sum_Dw_D\bigl(\tfrac DN\bigr)$; and simulate the random multiplicative series with independent signs $\varepsilon_q$. Deliverables: the residuals of the prediction, the empirical law of $K_N-2h$ against the simulated law (quantile plot and Kolmogorov Smirnov statistic), and the levels with the most extreme values together with their signs $\varepsilon_q$ for $q\le100$. Why it matters: agreement would support reducing Open Problem~\ref{op:kemeny} to the joint equidistribution of Legendre symbols, and the extreme levels test whether $K_N-2h$ is bounded.
\end{compproblem}
